The prover-verifier game framework has been shown to be useful in training legible language models. However, could we translate the same formalization for AI Control? In this section, we develop a control-robust variant that preserves the theoretical guarantees of the original while addressing its vulnerabilities to strategic deception.

\section{Preliminaries}
We work with measurable spaces $(\mathcal{X}, \mathcal{B}_X)$ and $(\mathcal{Z}, \mathcal{B}_Z)$ representing problems and solutions respectively, with a fixed distribution $P_X$ over problems. For each problem $x \in \mathcal{X}$, we assume access to either a ground-truth oracle $y: \mathcal{X} \to \mathcal{Y}$ or a verification function $C: \mathcal{X} \times \mathcal{Z} \to \{0,1\}$ that determines whether solution $z$ satisfies the specification for problem $x$. We denote by $\mathbbm{1}_{\text{correct}}(x,z) := C(x,z)$ the correctness indicator.
\\
The game involves two entities: a \emph{prover} and a \emph{verifier}. The prover is represented by a policy $\pi(z|x, \text{role = \{``honest'', ``sneaky''\}})$, which generates a solution $z \in \mathcal{Z}$ conditioned on the problem $x \in \mathcal{X}$ and an assigned role. The role determines the prover's immediate objective. Specifically, the role could be:
\begin{itemize}
    \item \textbf{Honest.} The prover, denoted $h(x) = \pi(x, \text{role=``honest''})$, aims to produce a correct solution that is deemed convincing by the verifier.
    \item \textbf{Sneaky.} The prover, denoted $s(x) = \pi(x, \text{role=``sneaky''})$, aims to produce an incorrect solution that nevertheless convinces the verifier\footnote{For our purpose, it is irrelevant whether the prover policy is unique and role conditioned --namely, the parameter space is shared among the two modes-- or whether the two roles make use of separate copies of the same initial policy. While shared policies are common in symmetric games \citep{silver2017alphazero}, if the game itself is fundamentally asymmetric (e.g., one player is a ``defender'' with different rules/objectives than an ``attacker'', or different starting conditions that aren't just a color swap), then separate networks might be essential because the optimal strategies are inherently different.}.
\end{itemize}
Both $h(x)$ and $s(x)$ are assumed to be drawn from the same underlying model class, denoted $\mathcal{H}$. Theoretically, $\mathcal{H}$ is considered to be the set of all deterministic mappings from $\mathcal{X}$ to $\mathcal{Z}$, although in practice we model $\mathcal{H}$ as a Large Language Model (LLM).
\\
The verifier, denoted as $v_\phi$, is similarly considered to be the set of all deterministic mappings $\mathcal{V}$, although it is still modeled as an LLM in our framework. The verifier's task is to determine the correctness of the prover's solution. Formally, $v_\phi: \mathcal{X} \times \mathcal{Z} \rightarrow [0,1]$ outputs a \emph{convincingness score} representing the verifier's belief that the solution $z$ is correct for problem $x$. This can be interpreted as $v_\phi(x, z) = P(C(x, z) = 1)$, the verifier's estimate of the probability that the solution is valid. Finally, all random variables are defined on a common probability space $(\Omega, \mathcal{F}, \mathbb{P})$.
\\
A fundamental requirement for this game, following \citet{kirchner2024prover}, is \emph{checkability}—the verifier's capacity to independently assess the correctness of prover-generated solutions. Checkability comprises two essential properties:

\begin{enumerate}
    \item \textbf{Completeness:} The verifier correctly accepts all valid solutions from the honest prover. Formally:
    \[\mathbbm{1}_{\text{correct}}(x, \pi(x)) = v(x, \pi(x)) = 1, \qquad \forall x \in \mathcal{X}\]

    \item \textbf{Soundness:} The verifier correctly rejects all invalid solutions, regardless of which prover in the model class generated them. Formally:
    \[\mathbbm{1}_{\text{correct}}(x, \pi'(x)) = 0 \implies v(x, \pi'(x)) = 0, \qquad \forall \pi' \in \mathcal{H}, \forall x \in \mathcal{X}\]
\end{enumerate}

These definitions draw inspiration from the literature on Probably-Approximately Correct (PAC) verification and interactive proof systems \citep{goldwasser2021interactive, pmlr-v195-mutreja23a}. However, in informal domains, perfect completeness and soundness are unrealistic expectations for both models and humans. Therefore, we study relaxed versions of these properties: rather than perfect completeness, we expect the honest prover to achieve a high pass rate and the verifier to demonstrate high recall on honest solutions ($\varepsilon$-completeness); instead of perfect soundness, we expect the verifier to maintain a low false positive rate when evaluating adversarial inputs ($\varepsilon$-soundness).

\section{Simultaneous Play for Legibility}
In order to formalize the original PVG setup and justify our modifications to it, let us establish the baseline components and provide some definitions. \citet{kirchner2024prover} proposes to train \textit{legible} prover and sound verifier using variants of the prover-verifier game \citep{anil2021learning}. They formulate the prover-verifier game as a simultaneous-move game between a role-conditioned prover $\pi_\theta$ and a verifier $v_\phi$. The prover, parameterized by $\theta$, generates solutions conditioned on its assigned role:

$$z_h = \pi_\theta(x, \text{helpful}), \quad z_s = \pi_\theta(x, \text{sneaky})$$

The verifier $v_\phi: \mathcal{X} \times \mathcal{Z} \to [-B, B]$ outputs an unconstrained convincingness score. The training proceeds via alternating optimization of two objectives. The verifier minimizes a classification loss:

$$L_V(\phi; \theta) = \frac{1}{2}\mathbb{E}_x\left[\ell(v_\phi(x, z_h), \mathbbm{1}_{\text{correct}}(x, z_h))\right] + \frac{1}{2}\mathbb{E}_x\left[\ell(v_\phi(x, z_s), \mathbbm{1}_{\text{correct}}(x, z_s))\right]$$

where $\ell$ is typically the cross-entropy loss or the Bradley-Terry loss. The prover maximizes a reward that encourages role-aligned behavior:

$$R_P(\theta; \phi) = \frac{1}{2}\mathbb{E}_x[r(v_\phi(x, z_h), \mathbbm{1}_{\text{correct}}(x, z_h))] + \frac{1}{2}\mathbb{E}_x[r(v_\phi(x, z_s), 1 - \mathbbm{1}_{\text{correct}}(x, z_s))]$$

The reward function $r(v, q)$ satisfies two critical properties:

\[\begin{aligned} \forall v_1, v_2 \in [0, 1] &: r(v_1, 1) > r(v_2, 0), \\ \forall v_1, v_2 \in [0, 1], v_1 > v_2 &: r(v_1, 1) > r(v_2, 1). \end{aligned}\]

\noindent where (i) states that role-aligned solutions always receive higher reward than misaligned ones, and (ii) states that among role-aligned solutions, higher verifier scores yield higher rewards.
\\
This formulation achieves notable empirical successes in the pursuit of legibility. Human evaluators rate helpful solutions as increasingly legible over training rounds, and the verifier demonstrates improved robustness against held-out adversarial examples. These properties suffice when the primary concern is unintentional errors, algebraic mistakes, logical non-sequiturs, or computational slips.

\subsection{Limitations for Strategic Deception}
When facing intentionally deceptive agents in control scenarios, three pathologies emerge that fundamentally limit the applicability of simultaneous-play PVG:

\begin{itemize}
    \item \textbf{Style Leakage and Spurious Correlations.} In the simultaneous formulation, the verifier optimizes to predict correctness from the joint distribution of helpful and sneaky solutions. Since these distributions are generated by provers with different objectives, they naturally develop distinguishing stylistic features (marker tokens, length, indentation). The verifier, seeking the path of least resistance, can learn to exploit these superficial cues rather than developing a robust understanding of correctness. This is a weakness that \citet{kirchner2024prover} highlight as well in \S C of their paper.
    % Formally, there exists a Lipschitz function $g: \mathcal{Z} \to \mathbb{R}$ such that $\mathbb{E}[g(z_h) - g(z_s)] > 0$ while the verifier's gradient $\mathbb{E}[\partial_v L_V] \approx 0$.
    The verifier has found a spurious feature that correlates with the helpful/sneaky distinction but not with actual correctness. This manifests empirically when verifiers trained on one pair of provers fail to generalize to new adversaries with different stylistic signatures. For this reason, the simultaneous game can converge to a \textbf{weak equilibria}, resembling (implicit) \textbf{collusive dynamics} between the honest prover and the verifier.
    \item \textbf{Backdoor Blindness and Distribution Mismatch.} Most critically for control applications, the original formulation never exposes the verifier to the specific trigger inputs on which backdoored solutions misbehave. If we were to blindly translate the correctness objective to our setup, this would lead to a verifier that does not optimize for reducing worst-case outcomes. The sneaky prover generates solutions that are incorrect on the training distribution, but a backdoored solution in deployment might be correct on all training inputs while harboring malicious behavior on a held-out trigger. Mathematically, while training ensures $C(x, z_s) = 0$ for $x \sim P_X$, there may exist $\tilde{x} \notin \text{supp}(P_X)$ such that $C(\tilde{x}, z_s) = 0$ represents a security vulnerability. The training objective provides no incentive for the verifier to detect such backdoors, as they manifest as correct solutions during training.\footnote{In \citet{kirchner2024prover}, sneaking in the training distribution is equivalent to sneaking in the test distribution, since being sneaky amounts to inserting a deliberate error anywhere in the reasoning trace. Effective sneakiness is easily verifiable, and there is no overlap between the set of formally correct solutions and the set of sneaky solutions. In code examples, this is not the case.}
\end{itemize}

\section{Sequential Play with Distribution Discrimination}

To address these limitations, we introduce three key modifications that transform the game from a legibility-focused framework to one suitable for adversarial control scenarios.

\paragraph{Sequential Stackelberg Structure.} We replace simultaneous play with a three-stage sequential game:
\begin{enumerate}
    \item The verifier $v_\phi$ commits to its evaluation strategy;
    \item The helpful prover $h(x)$ generates solution $z_h = h(x)$;
    \item The sneaky prover $s(x)$ observes $z_h$ and generates $z_s = s(x, z_h)$
\end{enumerate}

This sequential structure changes the game dynamics. The sneaky prover, observing the helpful solution, can copy any stylistic signals but must still satisfy its objective of introducing errors or backdoors. This forces genuine differentiation based on correctness rather than style.

\paragraph{Distribution Discrimination Objective.} Rather than training the verifier to predict correctness, we train it to identify which role distribution generated a solution. The verifier now aims to distinguish $P_{\text{honest}}$ from $P_{\text{sneaky}}$ rather than correct from incorrect solutions. Formally, $v_\phi$ learns $P(\text{role}=H\mid x,z)$. Let
\[
\mathcal D_H=\bigl\{(x,z):z=h(x)\bigr\},\qquad
\mathcal D_S=\bigl\{(x,z):z=s(x)\bigr\}.
\]

The Kirchner verifier minimises the \textit{classification risk}

\[
R^\text{leg}(v)=\Pr_{(x,z)\sim\mathcal D_H\cup\mathcal D_S}
               \Bigl[v(x,z)\neq \mathbbm{1}_{\text{correct}}(x,z)\Bigr],
\]

whereas our verifier minimises

\[
R^\text{ctl}(v)=\Pr_{(x,z)\sim\mathcal D_H\cup\mathcal D_S}
                \Bigl[v(x,z)\neq\mathbf 1[(x,z)\in\mathcal D_H]\Bigr].
\]

While $\mathcal D_H$ may contains some incorrect samples (when $h$ errs) and $\mathcal D_S$ contains some correct samples (when the attack passes all unit test instances), the two risks may disagree on a non-empty set of points.  The solutions of the two ERM problems therefore differ unless $h$ is perfect, which is, in the limit, the outcome that $h$ aims to maximize.
\\
We implement this via a modified Bradley-Terry objective:
$$L_V(\phi) = -\mathbb{E}_{x,z_h,z_s}\left[\log \sigma(v_\phi(x, z_h) - v_\phi(x, z_s))\right] + \lambda ||v_\phi||^2$$
\noindent~By learning to recognize the generative process rather than just the output property, the verifier is exposed to an even stronger adversarial pressure by the sneaky prover, which in turn helps it develop more robust features that transfer across different instantiations of sneaky behavior.
\\
The helpful reward remains straightforward:
$$R_H = v_\phi(x, z_h) \cdot \mathbbm{1}_{\text{correct}}(x, z_h)$$
\\
while the sneaky reward incorporates backdoor success:
$$R_S = v_\phi(x, z_s) \cdot \text{success}(x, z_s)$$
\\
Regarding the nature of the game, our prover-verifier game is still formulated as a general-sum game, not a zero-sum game: it is possible for both players to improve their utilities simultaneously.

\section{Theoretical Guarantees under Sequential Play}
Despite these substantial modifications, we can establish that the sequential game preserves the key theoretical properties of the original formulation.
\begin{definition}[Sequential Controllability]
A pair $(v, \pi_H)$ is $(\alpha, \beta)$-sequentially controllable if:
\begin{enumerate}
    \item ($\alpha$-Completeness) $\mathbb{E}_x[\mathbbm{1}_{\text{correct}}(x, \pi_H(x))] \geq 1 - \alpha$ and $\mathbb{P}[v(x, \pi_H(x)) > v(x, \pi_S(x, \pi_H(x)))] \geq 1 - \alpha$
    \item ($\beta$-Soundness) For all measurable $\pi': \mathcal{X} \times \mathcal{Z} \to \mathcal{Z}$,
   $$\mathbb{P}[\mathbbm{1}_{\text{correct}}(x, \pi'(x, \pi_H(x))) = 0 \wedge v(x, \pi_H(x)) \leq v(x, \pi'(x, \pi_H(x)))] \leq \beta$$
\end{enumerate}
\end{definition}

\begin{theorem}[Equivalence]
Under standard regularity conditions, an $(\varepsilon_V, \varepsilon_H, \varepsilon_S)$-approximate Stackelberg equilibrium of the sequential game is equivalent to an $(\alpha, \beta)$-sequentially controllable configuration with explicit moduli relating the parameters.
\end{theorem}
The proof, detailed in Appendix A, establishes that the sequential structure does not compromise the convergence guarantees while providing stronger robustness properties. The key insight is that the verifier's ability to distinguish distributions translates directly to detecting deceptive behavior, even from more capable future adversaries.

\section{Addressing the Original Pathologies}
The sequential reformulation attempts to address the limitations of the original framework:

\paragraph{Style Leakage Mitigation.} Sequentiality forces stylistic alignment between honest and sneaky outputs.This eliminates the simple strategy of developing distinguishing stylistic markers, or, without ascribing intentionality, any strategy relying on spurious stylistic correlations. The verifier must learn features genuinely indicative of the generative process rather than superficial correlates.

\paragraph{Avoiding Weak Equilibria.} The sequential best-response structure ensures the honest prover cannot simply coordinate with the verifier on a shared signaling strategy. After observing $z_h$, the sneaky prover could copy any coordinated signal, but must still introduce actual errors or backdoors to achieve its reward. This forces the verifier to develop genuinely discriminative capabilities.

\paragraph{Backdoor Detection.} By explicitly incorporating trigger demonstrations in the training process, the verifier learns to recognize the patterns of backdoored solutions. The perturbation testing ensures robustness—a backdoor that only works on precisely $\tilde{x}$ but fails on slight variants receives low reward, encouraging the sneaky prover to create more robust backdoors that the verifier must learn to detect.
\\
This reformulation transforms the prover-verifier game from an elegant framework for improving legibility into a practical tool for adversarial robustness in control scenarios. The theoretical guarantees ensure convergence to meaningful equilibria, while the structural modifications address the specific challenges of defending against strategic deception.